\notepage{63}{Independent orbital and spin degrees of freedom}
Suppose the Hamiltonian separates into independent spatial and spin
parts,
\[
 H=H_0\otimes I_s+I_r\otimes H_s,
\]
often abbreviated to $H=H_0+H_s$. Here $H_0$ contains only spatial
operators, and $H_s$ contains only spin operators.

By contrast, an interaction $H_{so}$ involving both sectors cannot
generally be written as such a sum. A familiar example is
\[
 H=H_0+H_s+H_{so},\qquad H_{so}=\xi\,\mathbf L\cdot\mathbf S,
\]
with a coupling coefficient $\xi$ (which can depend on radius in
a central-field model).

\smallheading{Separated time evolution}
For time-independent independent parts, a product initial state stays
factorized:
\[
 \ket{\Psi(t)}=\ket{\psi_0(t)}\otimes\ket{\chi(t)},
\]
where
\[
 \ket{\psi_0(t)}=\ee^{-\ii H_0t/\hbar}\ket{\psi_0(0)},\qquad
 \ket{\chi(t)}=\ee^{-\ii H_st/\hbar}\ket{\chi(0)}.
\]
\editorial{The spin propagator acts on the initial spin ket
$\ket{\chi(0)}$. A factorized Hamiltonian alone does not imply that
every initial state is a product; an initially entangled state
still evolves by the product of the two propagators.}

\smallheading{The spin factor alone}
For spin one half,
\[
 \ket{\chi}=\alpha\ket{+}+\beta\ket{-}
 \quad\longleftrightarrow\quad
 \begin{bmatrix}\alpha\\\beta\end{bmatrix},
 \qquad
 \ket{+}\longleftrightarrow\begin{bmatrix}1\\0\end{bmatrix},\quad
 \ket{-}\longleftrightarrow\begin{bmatrix}0\\1\end{bmatrix},
\]
with $\alpha=\braket{+}{\chi}$ and $\beta=\braket{-}{\chi}$.
Inserting the spin completeness relation gives
\begin{align*}
 1=\braket{\chi}{\chi}
 &=\sum_{m_s=\pm1/2}\braket{\chi}{m_s}\braket{m_s}{\chi}\\
 &=\abs{\braket{+}{\chi}}^2+\abs{\braket{-}{\chi}}^2
 =\abs\alpha^2+\abs\beta^2.
\end{align*}
